\chapter{Experimental Methodology}
\label{ch:experiments}
The experiments address detection, numerical computation, and adaptation
through separate evidence bundles. They do not form one interchangeable
collection of scores. This chapter states the input roles, settings,
denominators, and measurement boundaries before presenting their results.
In particular, the historical DM-only experiments, retrospective scoring
selection, matched-duration study, and closing final-partition replay are
identified separately.

\section{Datasets, devices, and label provenance}
\label{sec:experimental-data}

\subsection{Two purposes for the available traffic}

D-Link Camera~5 supplies the count workload connected to the inherited
DM IoT study \cite{alhajbaddar2025iot}. The underlying D-Link traces were
published by Roy Chowdhury et al. \cite{chowdhury2021data}. The project
selects twelve captures from October 2020 for fitting, calibration, and
later numerical/scoring workloads. Attack labels are unavailable for
these selected captures. D-Link alert counts can therefore describe
detector output but cannot estimate attack recall, precision, or a false
positive rate.

The labelled detection experiments use the UNSW consumer-IoT attack
collection associated with Hamza et al. \cite{hamza2019,unswattackdata}. Mixed-device
captures are filtered separately for Chromecast and Samsung Camera.
Chromecast uses MAC address \path{f4:f5:d8:8f:0a:3c}; Samsung uses
\path{00:16:6c:ab:6b:88}. These views are not the UNSW-NB15 dataset and
are not N-BaIoT. Keeping their identity explicit avoids attributing a
published result from another benchmark to the present traces.

The source files remain separate from processed windows and labels.
Prepared source directories can use symbolic links, avoiding a second
copy of the capture bytes. Device filtering changes the analytical view,
not the publisher's files. Model artifacts record the selected device
and captures so that a directory containing many sources is not mistaken
for the set actually used in an experiment.

\begin{table}[tbp]
\centering\small
\begin{tabular}{@{}>{\raggedright\arraybackslash}p{.19\textwidth}
 >{\raggedright\arraybackslash}p{.14\textwidth}
 >{\raggedright\arraybackslash}p{.26\textwidth}rrr@{}}
\toprule
\textbf{Device} & \textbf{Role} & \textbf{Selected date range} &
\textbf{Windows} & \textbf{Attack} & \textbf{Silent} \\
\midrule
\TableRows{generated/partitions}
\bottomrule
\end{tabular}
\caption{Reference ten-minute observations, generated from the configured
processed captures and label files. Date ranges give the first and last
selected nominal capture IDs and do not assert continuous coverage. D-Link
has no attack labels. Silent windows are observed zero vectors and remain
part of the window count.}
\label{tab:partitions}
\end{table}

\subsection{Coverage, exclusions, and silent observations}

Table~\ref{tab:partitions} gives the stored reference composition. The
D-Link development captures are October~16 and~19, rather than every day
between those dates. Its fit set has 846 rows, its calibration 288,
development 270, and later numerical partition 288. Chromecast has 867
fit, 576 calibration, 474 development, and 432 final observations.
Samsung has 432 fit, 81 calibration, 556 development, and 576 final
observations. These are counts of retained observations, not assumptions
that every nominal day contains exactly 144 ten-minute intervals.

Chromecast's October~24 capture is excluded because its final PCAPNG
block is truncated. The preparation record documents that another
download did not repair the publisher-side source. Six annotation
intervals on that date consequently have no retained capture. October~21
contains only about seven hours, producing 42 windows. Samsung May~31
contains 81 windows and June~1 contains 124. Metrics are evaluated on
matched observed windows rather than completing these dates with
unobserved data.

The stored Chromecast development set contains 187 silent windows;
Samsung contains 30 silent development windows and 19 silent final
windows. These observations are distinct from missing coverage and are
retained in the denominators. Their presence helps explain why a large
number of true negatives need not imply broad coverage of active device
behaviour. No missing-state rows are stored in the selected labelled
collections, but that does not prove the original recording was
continuous outside the observed boundaries.

An additional chronology check for this thesis found 846 D-Link fitting
rows but only 702 distinct UTC window bounds: 144 bounds are repeated
across selected captures. The saved model is preserved because the
D-Link comparison concerns numerical implementations on the same retained
matrix. The rows must not be described as 846 independent time intervals.
Resolving that overlap and refitting would define a different workload;
it is not silently substituted for the published project evidence.

\subsection{Annotation preparation and alignment}

The original annotation bytes are fingerprinted separately for each
selected device. Chromecast's official file contains 27 intervals, of
which 21 match retained captures; Samsung's 42 intervals all match.
The preparation record retains the unmodified annotation identity and
derives per-window labels separately from the processed count files.
Labels are not appended to the training matrices or consumed by the
detector's admission rule.

For annotated interval $[a_s,a_e)$ and observation $[w_s,w_e)$, the
window is attack-positive if
\begin{equation}
 a_s<w_e\quad\text{and}\quad a_e>w_s.
\label{eq:label-overlap}
\end{equation}
Any positive overlap is sufficient, regardless of the fraction of the
window occupied by the attack. When intervals overlap, their union
determines overlap seconds so that a window cannot contain more labelled
attack duration than its own length. Exact UTC bounds, rather than only
the filename date, determine this alignment.

The prepared Chromecast labels contain 38 positive windows across the
retained data: 32 development and six final. Samsung has 82: 70 development
and twelve final. An attack interval can create more than one positive
window and can overlap another annotation. The number of positive windows
is therefore neither an attack count nor an estimate of the duration of
malicious packet transmission. The thesis reports both window and
annotation-interval outcomes.

\section{Chronological partitions and leakage prevention}
\label{sec:experimental-chronology}

\subsection{Reference fit, calibration, and evaluation roles}

Complete captures are assigned to roles. The fit period precedes benign
calibration, followed by development and the configured final period.
The labelled fit and calibration partitions contain no attack-positive
windows under the supplied annotation scheme. That supports their stated
role as presumed normal data; incomplete annotations could still leave
unrecognised activity in the reference. The assumption is recorded rather
than treated as a universal guarantee of benignness.

Fitting reads the fit role, and calibration reads its distinct role.
Model and threshold artifacts preserve source identities and time bounds.
The evaluation joins events and labels on capture context and exact
window bounds, rejecting duplicate or incompatible records and reporting
unmatched inputs. The dedicated closing replay checks actual chronological
overlap and complete label coverage as well as nominal capture membership.

This protocol avoids randomly mixing earlier and later windows in a
single train/test split. It does not remove temporal dependence or produce
multiple independent histories. Adjacent observations can be correlated,
and a limited sequence of captures can favour particular update times.
The separate temporal branch studies rolling folds and episode-level
summaries; its outputs are not presented as results from the retained
detector. This boundary follows the broader concern with time and
distributional bias in security evaluation \cite{pendlebury2019}.

\subsection{Development selection and prior exposure}

Development labels were used to compare scoring choices, adaptation
schedules, and window durations. The learning calculations still use
normal fitting/calibration data, but the research process is not blind
to labelled outcomes. The broad per-category-tail experiments also
inspected later-partition labels; their own record states that subsequent
comparisons are retrospective. A later replay cannot undo that exposure.

The term \emph{final partition} in this thesis identifies the saved role
and the closing evaluation bundle. It does not claim an independently
untouched test set. The closing run fixes the selected ten-minute rule
and starting pairs before replay, does not select settings from its
outputs, and saves the complete records. Its value is reproducibility and
coverage of that declared later sequence. Independent external
generalisation remains unmeasured.

This distinction is necessary even when the same settings are used on
both labelled devices. Shared settings reduce device-specific tuning
but do not erase the design choices informed by their data. Arp et al.'s
discussion of data snooping and experimental pitfalls motivates explicit
reporting of these choices \cite{arp2022}. The comparison tables label
historical development, retrospective pooled, duration-specific later
audit, and closing final results separately.

\section{Static and adaptive configurations}
\label{sec:experimental-configurations}

\subsection{Common reference settings}

The reference detector uses the four categories and a ten-minute duration,
the native LM backend at six requested digits, initial concentration~10,
likelihood-change tolerance $2\times10^{-9}$, and at most 100,000 fitting
iterations. Joint calibration uses the 0.01 lower quantile with at least
30 observed normal calibration windows. Every mode starts from a saved
normal-only model and its matching threshold, not from the final state
left by a previous mode.

The three modes are static, adaptive threshold, and periodic refit.
The conservative settings use a 288-window buffer, a 144-observation
attempt cadence, 144 accepted windows for a refit, confidence multiplier
one, and the 0.99 calibration-volume guard. Concentration may change by
at most a factor~1.25 per candidate, the largest category-probability
shift is limited to 0.05, and threshold movement is bounded by twice the
original IQR scale. These constants are retained engineering choices;
they were not estimated as optimal security bounds.

Historical DM-only comparisons retain the legacy threshold and refit
rules. The ordinary refit uses cadence/minimum 144/144; the selected
development variant uses 288/240. The shared 288-window buffer is not the
same parameter as the cadence. Retrospective DM-only versus combined
comparisons use 288/144 and include both policies. The thesis never
attributes their pooled improvement to the historical 288/240 run.

\subsection{Scoring comparisons and rejected alternatives}

The retained combined detector is compared with the earlier complete
DM-only score. Both fit the same distribution and use the same count
features; they differ in the calibrated evidence combined at inference
and in preservation of original detection references. The retrospective
pool comprises the two devices' development and later partitions,
2,038 labelled windows in total. Each mode and partition starts
independently. Counts are summed before pooled metrics are computed.

Per-packet normalisation, broad upper tails for every category,
separately allocated extreme-tail cutoffs, and joint rank-tail calibration
are retained as rejected choices. They are reported because they explain
selection of the small final rule, not because their best individual
scores can be assembled into a new detector after inspecting labels.
Additional packet fields, including byte, direction, and SYN information,
were not adopted without a common validated input across both devices.

No Poisson detector, N-BaIoT implementation, or Kitsune implementation was
rerun under this protocol. Their papers supply context in related work,
not a controlled performance baseline. Conclusions about efficiency and
accuracy are accordingly restricted to the comparisons actually made.

\subsection{Matched window-duration study}

The duration study compares 10, 20, 30, and 60 minutes on common
capture-local observed support. Complete observed hourly blocks define
the shared support, so every duration aggregates the same retained traffic.
Counts are summed, and annotations are reapplied at the new boundaries.
The study excludes unmatched capture edges and records the exclusions;
it does not join incomplete pieces from different captures into an hour.

Every duration receives a fresh normal fit and joint calibration. The
normal split is the same across durations within this study. To retain
at least thirty calibration observations at sixty minutes, the latest
whole normal fitting capture is moved to calibration where necessary.
Samsung therefore fits May~28--29 and calibrates May~30--31, producing
37 hourly calibration windows. Chromecast retains its normal split and
has 96. The rule depends on normal sample support, not attack labels.
It makes the study's Samsung ten-minute model different from its saved
reference model; this difference is stated rather than hidden.

Adaptive settings scale to a 48-observed-hour buffer, 24-hour attempt
cadence, and nominal 24 accepted hours before refitting. The unchanged
thirty-sample conservative guard implies an effective minimum of thirty
accepted hours at sixty-minute resolution. The guard is not loosened to
produce an update. Each evaluation role and mode again starts independently.
Forty-eight combinations cover two devices, four durations, two roles,
and three modes.

\section{LM and SciPy comparison protocol}
\label{sec:experimental-backends}

\subsection{Production row-wise workload}

The retained D-Link production benchmark uses an 846-by-4 fitting matrix,
288 calibration windows, and 270 development windows. The two backends
receive the same observations, estimator settings, initialisation rule,
and complete-score convention. Temporary artifacts keep the comparison
from replacing accepted starting pairs. Agreement measures include
alpha differences, likelihood and threshold differences, score
differences, and decision mismatches.

Complete-fit timing excludes source loading and artifact saving. Five
fit measurements provide the median. Kernel timing uses 200 calls per
batch and single-window timing 1,000 calls per batch, with five batches
and warm-up. The representative count vector is $(44,112,88,166)$.
This separates a many-row training kernel from a one-row scoring call,
whose Python/native and coefficient overheads have different relative
costs.

The production record reports Python~3.12.14, NumPy~2.5.1, and
SciPy~1.18.0 on Linux x86-64. A distinct combined-detector backend-check
record covers all three devices and their development observations.
Its ratios are reported separately, rather than substituted into the
D-Link timing table. Source and environment differences remain attached
to each result.

\subsection{Aggregate architecture and numerical stress probes}

A separate experiment follows the supplied study's aggregate convergence
architecture: seven D-Link days, four sampling fractions
(10/15/20/25~percent), three fixed repeats, row-wise digamma updates,
and a pooled-vector stopping likelihood. Both backends use identical
samples. The resulting 84 paired fits answer a controlled architectural
comparison on local data. Their convergence objective differs from the
production row-wise objective, so their timing is not used to justify
changing that production contract.

The source publication's exact random samples cannot be recovered, and
local daily row counts differ. The experiment is therefore not a literal
reproduction of every published table. It evaluates the same broad
architecture with controlled backend inputs. Native setup optimisation
and high-concentration probes are likewise treated as their recorded
experiments rather than new runs conducted while writing this thesis.

The ordinary workload can show close agreement even when extreme
concentrations expose cancellation. High-precision probes provide
additional evidence about that regime, distinct from attack detection.
No conclusion that smaller numerical error improves recall is drawn
without observing a changed decision on labelled traffic.

\section{Detection metrics and attack episodes}
\label{sec:experimental-metrics}

\subsection{Window classification}

Let TP, FP, FN, and TN count matched windows with the usual true/false
positive and negative outcomes. The principal rates are
\begin{equation}
 \mathrm{Precision}=\frac{\mathrm{TP}}{\mathrm{TP}+\mathrm{FP}},\qquad
 \mathrm{Recall}=\frac{\mathrm{TP}}{\mathrm{TP}+\mathrm{FN}},\qquad
 \mathrm{FPR}=\frac{\mathrm{FP}}{\mathrm{FP}+\mathrm{TN}}.
\label{eq:classification-rates}
\end{equation}
For defined precision and recall with positive sum,
\[
 \mathrm{F1}=\frac{2\,\mathrm{Precision}\,\mathrm{Recall}}
                  {\mathrm{Precision}+\mathrm{Recall}}.
\]
Accuracy is also available, but can be dominated by the many benign
windows. Confusion counts accompany the rates, consistent with the
class-imbalance concerns described by Saito and Rehmsmeier
\cite{saito2015}.

The evaluator stores undefined quantities as JSON null. In particular,
zero denominators do not become perfect or zero performance by default;
the current F1 implementation also returns null if precision and recall
sum to zero. All principal tables here have positive true-positive counts,
so their displayed F1 values are defined. Pooled F1 is calculated from
summed confusion counts, not the arithmetic mean of device F1 values.

ROC--AUC describes score ordering where the score is available, with
orientation chosen so more anomalous observations are more positive.
Device-wise results have a clear source boundary. A pooled AUC across
device-calibrated margins is only descriptive because those margins need
not have a common ranking interpretation. The primary closing result is
the binary confusion table. Later count-first scoring can preserve that
table while changing severity and AUC.

\subsection{Annotated intervals, coverage, and response time}

An interval is counted as detected if at least one alerting window
overlaps that official annotation under Eq.~\eqref{eq:label-overlap}.
Only fully observed intervals enter the common denominator in the
duration study. Partially observed exclusions are reported. These
annotation intervals are not merged into a different incident definition,
and they are not automatically equivalent to the range metrics of Tatbul
et al. \cite{tatbul2018}.

The first operational alert time is the end of the first overlapping
alert window. If that time is $t_a$ and attack onset is $a_s$, delay is
$t_a-a_s$. A window that overlaps an attack may close after its annotated
end. The number alerted by attack end is therefore reported separately
from the number eventually detected. Median delay includes detected
intervals only; missed attacks are not silently assigned zero delay.

In the duration comparison, window F1 changes its unit as duration
changes. A second summary projects each duration's decision onto its
constituent ten-minute bins and compares those bins with their labels.
This common-grid F1 assesses localisation on a fixed reference grid.
It does not move the operational alert earlier: a thirty-minute decision
still becomes available at the thirty-minute window's close.

False-alert frequency per 24 observed hours is
$24\,\mathrm{FP}/H$, where $H$ is the retained observation time in hours.
It is distinct from FPR in Eq.~\eqref{eq:classification-rates}, whose
denominator contains benign-labelled windows. Reporting both separates
the relative classification error from a workload-normalised alert count.

\subsection{Finite calibration support and tail resolution}
\label{sec:tail-resolution}

The nominal one-percent quantile is the same across reference models,
but its empirical support is not. With linear interpolation and
ordered calibration scores $J_{(1)}\leq\cdots\leq J_{(n)}$, the
one-based interpolation position for probability $q$ is
\begin{equation}
 h=1+(n-1)q.
\label{eq:quantile-position}
\end{equation}
Writing $j=\lfloor h\rfloor$ and $\lambda=h-j$, the quantile is
$(1-\lambda)J_{(j)}+\lambda J_{(j+1)}$, with the usual endpoint
convention. This is the linear sample-quantile convention described
by Hyndman and Fan \cite{hyndman1996}. The position is not a count
of independent tail events.

\begin{table}[tbp]
\centering\small
\begin{tabular}{@{}lrr@{}}
\toprule
Calibration setting & Observations $n$ & Position $h$ at $q=0.01$ \\
\midrule
\TableRows{generated/quantile_ranks}
\bottomrule
\end{tabular}
\caption{Resolution of the empirical one-percent calibration tail.
Reference sample counts come from saved thresholds; the hourly sample count
comes from the matched-duration record. The minimum is a declared guard.}
\label{tab:quantile-ranks}
\end{table}

Table~\ref{tab:quantile-ranks} makes the difference concrete. The
reference Samsung calibration has 81 observations, so the selected
tail lies between the smallest and second-smallest values. Its
37-observation hourly calibration also depends on the first two
values, with a different interpolation weight. The 30-observation
minimum prevents a degenerate fit to only a handful of windows,
but does not turn those first two values into a reliable estimate
of an arbitrary future population tail.

Even the larger Chromecast calibration remains finite and temporally
structured. Repeated routine activity may yield many similar
observations, and adjacent windows need not provide independent
evidence about rare normal regimes. More recorded windows can
improve empirical resolution while still omitting a legitimate
operating condition. The evaluation therefore keeps the reference
period and traffic assumptions explicit rather than interpreting
sample count as complete benign coverage.

The practical implications extend to the joint rule. One unusually
low standardised feature can set the joint minimum and influence
the common quantile. This may leave another feature with a very
permissive implied cutoff, even though all features were scaled
by fitting medians and IQRs. Chapter~\ref{ch:results} consequently
reports the resulting count cutoffs as well as the nominal quantile.
The scale transformation and the sample-tail estimate solve
different problems.

The same limitation applies when accepted runtime observations
recalibrate a threshold. Those observations are selected by the
admission rule, not drawn independently from the full future benign
distribution. Recomputing a one-percent quantile on that buffer
does not guarantee a one-percent false positive rate. The candidate
history and observed errors must determine whether an update
helped; passing a minimum-size guard is insufficient evidence.

No confidence interval for a population false-alert rate is inferred
from Eq.~\eqref{eq:quantile-position}. Such a statement would need
assumptions about sampling and dependence that the available short
capture histories do not establish. The experiment instead reports
the exact calibration counts, fixed rule, and later confusion
matrices. This preserves the reproducible operating point while
keeping its statistical interpretation appropriately limited.

\section{Simulation environment and resource measurements}
\label{sec:experimental-resources}

The resource workload is a synthetic five-record sequence containing
ordinary traffic, an observed silent window, a missing window, and
alerting counts. Advancing timestamps permits accelerated repetitions
without violating chronology. Trials contain 2,000 and 20,000 supplied
records; four out of every five produce scored events. They begin from
isolated copies of the prepared conservative D-Link state, with the
288/144 learning schedule and unchanged movement/admission guards.

The execution environment is a Linux/amd64 Docker container with one CPU
and 512~MiB memory. The recorded host is an AMD Ryzen~7~9700X running
AlmaLinux~10.2; Docker is version~29.7.2. The measured image is Debian~13
with Python~3.11.17. This environment differs from the numerical
benchmark's Python~3.12 record. A CPU quota limits scheduling resources;
it does not emulate a particular embedded processor's clock, caches,
memory bus, or power consumption.

Timing begins when an input line is available and ends after its event
is flushed. It includes validation, scoring, attempted refits, candidate
handling, serialisation, and output. It excludes the time spent collecting
a ten-minute traffic interval. The synthetic producer contributes to total
replay throughput, while actual estimator calls are timed separately.
Rejected fits still contribute their full cost.

Process memory is measured through peak resource usage and sparse Linux
RSS samples. These counters have different accounting boundaries, so
both maxima are reported. State and log bytes, CPU time, image identity,
dependency versions, fixture/script hashes, and at most ten resource
samples are retained. The finite replay stores counters rather than an
unbounded event history in the detector process.

No update is forced by relaxing the guards. Thus a run with no promoted
candidate measures the rejection path and its cost. Controlled acceptance,
restart, rollback, signal, storage-failure, and retention tests provide
separate correctness evidence. A synthetic resource trial is not a
labelled security experiment or proof of indefinite bounded operation.

\section{Reproducibility and retained artifacts}
\label{sec:experimental-provenance}

The closing protocol is saved before its replay and identifies its
starting models, thresholds, configurations, code, processed inputs,
labels, and annotation identities. Each of the two devices is replayed
independently in each of three modes, retaining six event files and their
update attempts. The original reference pairs are not overwritten.
The recorded closing suite initially contains 601 passes; the later full
suite and the fresh thesis-side verification contain 643 passes with no
skips. Those counts describe different verification boundaries, not an
increase in the number of evaluated attacks.

The thesis presentation is generated from retained JSON and processed
records. The generator verifies arithmetic consistency of confusion
counts and rates, enumerates its source hashes, and checks the frozen
data/configuration inputs. It also detects three later source changes in
the probability, event-schema, and scoring modules. Their current hashes
match the subsequent scoring-order regression record, which reports
unchanged decisions and update outcomes on the same data. They are not
misreported as byte-identical to the earlier final source snapshot.

The paper-style aggregate experiment, production benchmark, combined-score
checks, and long resource replay each keep their own environment and
source record. The resource image was measured from a working tree based
on commit \texttt{8365870}, with the dirty status and package-source hash
recorded. A later commit name alone cannot identify that image precisely.
Likewise, the closing replay records its source boundary separately from
the present thesis-writing revision.

Following Sandve et al.'s provenance principle \cite{sandve2013}, the
reproducibility appendix maps the main claims to retained paths and
commands. It distinguishes work rerun for thesis verification from
historical measurements that were inspected. Raw capture preprocessing,
all fitting experiments, and the long resource replay were not repeated
merely to typeset the same values. No new external-device or physical
gateway validation is implied by a successful document build.
